\subsection{Original prescribed-class convexity on the actual native quotient chart}

\begin{definition}[The actual analytic projection of the original chart]
\label{projective:actual-admissible-analytic-projection}
\lean{EnvelopeProperty.ProjSubvariety.nativeAdmissibleAnalyticProjection}
\leanok
\uses{projective:actual-admissible-analytic-native-chart}
An actual native point of the original quotient-chart scheme over the trivial
coefficient valuation gives an original point of $X$ by applying the actual
original chart map and the inverse original analytic/native homeomorphism.
\end{definition}
\begin{proof}
The established original chart map preserves the actual coefficient structure.
It therefore maps the entire prescribed-base point space to that of the
original $X$, on which the inverse original homeomorphism is defined.
\end{proof}

\begin{lemma}[The actual projection equals the original chart inverse]
\label{projective:actual-admissible-projection-chart-inverse}
\lean{EnvelopeProperty.ProjSubvariety.nativeAdmissibleAnalyticProjection_eq_chartInv}
\leanok
\uses{projective:actual-admissible-analytic-projection,berkovich:actual-affine-paper-homeomorphism-inverse,projective:actual-admissible-analytic-native-chart}
For every actual original affine native point, this actual projection is the
original chart inverse of precisely its original affine seminorm point.
\end{lemma}
\begin{proof}
Apply injectivity of the original analytic/native homeomorphism. Its inverse
cancels the actual mapped point. The whole original chart comparison gives
the image of the original chart inverse. The canonical affine native point
of its original affine seminorm is the original native point itself, by the
two-sided affine homeomorphism identity. The entire mapped points agree.
\end{proof}

\begin{definition}[The actual original admissible weight]
\label{projective:actual-original-admissible-chart-weight}
\lean{EnvelopeProperty.ProjSubvariety.nativeAdmissibleChartWeight}
\leanok
\uses{projective:actual-admissible-analytic-projection}
For an original real function $\varphi$, an integer $m$, and an original
nonzero homogeneous $F$ of degree $d>0$, the actual chart weight is
$\varphi(w)-m^{-1}\log|F(w)|$, where $w$ is the actual preceding projection.
It is an extended-real function on the genuine original native quotient chart.
\end{definition}
\begin{proof}
Use the original normalized value of the same original polynomial at the true
projected point. The formula is defined for every original chart point; the
preceding chart identity also gives its nonvanishing denominator there.
\end{proof}

\begin{lemma}[Original chart convexity supplies genuine regular-unit convexity]
\label{projective:actual-original-admissible-chart-convexity}
\lean{EnvelopeProperty.ProjSubvariety.nativeAdmissibleChartWeight_isConvexPLOn}
\leanok
\uses{projective:actual-admissible-projection-chart-inverse,projective:actual-original-admissible-chart-weight,berkovich:actual-affine-convexity-from-paper}
If the original admissible coefficient-chart weight is convex PL on the whole
original affine chart, then the actual preceding weight is convex PL on the
whole actual native chart in the intrinsic definition using genuine units of
actual original scheme section rings.
\end{lemma}
\begin{proof}
Apply the original-affine to genuine-native regular-unit convexity theorem.
The actual projected point equals the original chart inverse by the preceding
lemma. Substitute that whole equality into the function, so its pullback is
literally the original admissible chart function. Its domain is the whole chart.
\end{proof}

\begin{corollary}[The prescribed rational class is retained]
\label{projective:actual-original-prescribed-class-chart-convexity}
\lean{EnvelopeProperty.ProjSubvariety.IsLocallyConvexAll.nativeAdmissibleChartWeight_isConvexPLOn}
\leanok
\uses{projective:actual-original-admissible-chart-convexity}
Original local $\mathcal O(a)$-convexity supplies the preceding genuine native
regular-unit convexity for every original admissible pair $m,d>0$ with the
exact original identity $d=ma$ and every original $F\notin I_X$ of degree $d$.
No nefness, curve-slope identity, or numerical positivity premise is supplied.
\end{corollary}
\begin{proof}
Specialize the unchanged original universally quantified local convexity
hypothesis to this original $m,d,F$, retaining its original equality $d=ma$.
Apply the preceding actual chart convexity theorem.
\end{proof}
